\documentclass[10pt,a4paper]{amsart}
\usepackage[margin=19mm]{geometry}
\usepackage{amsmath,amssymb,mathtools}
\usepackage[scr=rsfs,cal=euler]{mathalfa}
\usepackage{xfrac}
\usepackage[unicode,hidelinks,bookmarks=false]{hyperref}
\newcommand{\ie}{i.e.\ }
\newcommand{\cf}{cf.\ }
\newcommand{\resp}{respectively\ }
\newcommand{\Subsection}{\subsection}
\providecommand{\nobreakdash}{\nobreak\hbox{-}}
\setlength{\emergencystretch}{2em}
\multlinegap0pt
\usepackage{xcolor}
\usepackage{enumerate}
\usepackage[shortlabels]{enumitem}
\usepackage{amssymb}
\usepackage{tikz}
\newtheorem{lemma}{Lemma}
\newtheorem{proposition}{Proposition}
\newtheorem{claim}{Claim}
\def\Pb{\mathbb{P}}
\def\bE{\mathbb{E}}
\def\by{\boldsymbol{y}}
\def\cN{\mathcal{N}}
\def\eps{\varepsilon}
\title{\vspace{-0.9cm} Gaussian random graphs and Ramsey numbers}
\author{Zach Hunter, Aleksa Milojevi\'c, Benny Sudakov}
\date{Source: arXiv:2512.17718v2 (CC BY 4.0)}
\begin{document}
\maketitle

\noindent\emph{Source and licence.} \vspace{-0.9cm} Gaussian random graphs and Ramsey numbers. Version arXiv:2512.17718v2 (CC BY 4.0), distributed by arXiv under the Creative Commons Attribution 4.0 International licence, http://creativecommons.org/licenses/by/4.0/, as recorded in the arXiv licence metadata for this exact version and shown on the abstract page https://arxiv.org/abs/2512.17718v2. Full licence notice: \texttt{licenses/arxiv-2512.17718-CC-BY-4.0.txt}.\par\smallskip

\noindent\textit{Editorial scope.} Counted unit: the article's proposition prop:induction\_1 (source0.tex lines 501--510). The article's complete original proof of it is delivered with it (source0.tex lines 511--642, one \texttt{proof} environment of 132 lines). No mathematics is written, repaired or re-derived: the statement and the proof are the article's own lines, in the article's own notation, and no retained line is substituted or re-flowed.  Retained uncounted as the source-local context the proof needs: lines 314--316; lines 321--329; lines 350--355.  Nothing was omitted from any retained span, so the package carries no omission record.  No retained line contains a citation, so the wrapper carries no bibliography and no bibliography entry.  No retained line contains a figure environment, an image-inclusion command or drawing code, so no image is delivered and the package declares no asset dependency.  Editorial formatting only, in addition: an AMS \texttt{amsart} preamble that restates the article's own declarations and macros used by the retained lines, namely the environments \texttt{lemma}, \texttt{proposition}, \texttt{claim} on separate counters, together with the standard packages those lines need.  The retained environments print in this document as Lemma 1, Proposition 1, Claim 1 in order of appearance, whereas the article prints the same items with the numbering of its own full document.  The complete native source of the article as distributed in the arXiv e-print for this exact version is bundled unchanged as \texttt{sources/191340/source0.tex}.
\begin{equation}\label{eqn:log_concavity}
    \Phi(t+\eps)\leq \Phi(t)e^{\eps\phi(t)/\Phi(t)}.
\end{equation}

\begin{lemma}\label{lemma:truncated_expectation}
Let $b, \eps$ be real numbers and let $X\sim \cN(0, \frac{1}{d})$. Then, we have the following four formulas
{\smaller[0.3]
\begin{align*}
    &\bE\bigg[X\Big| X\leq \frac{b}{\sqrt{d}}\bigg]=-\frac{\phi(b)}{\Phi(b)\sqrt{d}} \quad \,\,\,\,\text{ and }\bE\bigg[X\Big| X\leq \frac{b+\eps}{\sqrt{d}}\bigg]=-\frac{\phi(b)}{\Phi(b)\sqrt{d}}+O\bigg(\frac{\eps}{\sqrt{d}}\bigg),\\
    &\bE\bigg[X\Big| X\geq \frac{b}{\sqrt{d}}\bigg]=\frac{\phi(b)}{(1-\Phi(b))\sqrt{d}}\text{ and } \bE\bigg[X\Big| X\geq \frac{b+\eps}{\sqrt{d}}\bigg]=\frac{\phi(b)}{(1-\Phi(b))\sqrt{d}}+O\bigg(\frac{\eps}{\sqrt{d}}\bigg).
\end{align*}
}
\end{lemma}

\begin{lemma}\label{lemma:exponential_moment_special_case}
Let $X_1, \dots, X_k$ be independent subgaussian random variables with variance proxy $1/d$, and let $\lambda\in (-\infty, \infty)$ such that $d\geq 4|\lambda| k$. If $S=\sum_{1\leq i<j\leq k}X_iX_j$, then
{\smaller[0.3]
\[\bE[e^{\lambda S}]\leq \exp\bigg( \lambda \bE[S] + \frac{\lambda^2k^2}{d}\sum_{j=1}^k (\bE X_j)^2 +\frac{4|\lambda| k}{d}\bigg).\]
}
\end{lemma}

\begin{proposition}\label{prop:induction_1}
Let $C>1, p\in (0, 1)$ be real numbers and $\ell, d=D^2\ell^2$ be large positive integers, where $D$ is sufficiently larger than $C$. Also, let $s, r$ be nonnegative integers such that $1\leq s+1\leq r\leq C\ell$. Suppose that the first $s$ columns of $M$ are fixed, and denote them by $M[s]=(M_1, \dots, M_{s})$. Then, the probabilities over the random choice of the remaining columns $M_{s+1}, \dots, M_r$ of the events $C_{s}\wedge B_r, I_{s}\wedge B_r$ are at most
{\smaller[0.3]
\begin{align}
    &\Pb\Big[I_s\wedge B_r\Big| M[s]\Big]\leq p^{\binom{r-s}{2}}\exp\bigg(\!-\!\frac{a\sqrt{d}}{p}\!\!\!\!\sum_{s<i< j\leq r}\!\!\!\!\langle \pi_s(\by_i), \pi_s(\by_j)\rangle \!-\!\frac{a^3}{p^3\sqrt{d}}\binom{r\!-\!s}{3}\!+\!O\Big(\frac{(r\!-\!s)^4}{d}\!+\!(r\!-\!s)\Big)\bigg)\label{eq:induction_bound_red}\\
    &\Pb\Big[C_s\wedge B_r\Big| M[s]\Big]\leq (1\!-\!p)^{\binom{r-s}{2}}\exp\bigg(\frac{a\sqrt{d}}{1\!-\!p}\!\!\sum_{s<i< j\leq r}\!\!\langle \pi_s(\by_i), \pi_s(\by_j)\rangle\!+\!\frac{a^3}{(1\!-\!p)^3\sqrt{d}}\binom{r\!-\!s}{3}  \label{eq:induction_bound_blue} \\
    &\hspace{11.0cm}\!+\!O\Big(\frac{(r\!-\!s)^4}{d}\!+\!(r\!-\!s)\Big)\bigg) \notag.
\end{align}
}
\end{proposition}

\begin{proof}
We argue by reverse induction on $s$. Note that for $s=r-1$, the statement is vacuous, since there the whole matrix is revealed and the product on the right hand side of Equation (\ref{eq:induction_bound_red},\ref{eq:induction_bound_blue}) is empty. Thus, we may use this as a basis of the induction. Let us first focus on bounding the probability of $C_s\wedge B_r$.

Suppose now that the statement holds for $s$ and let us prove it for $s-1$, where only the columns $M[s-1]=(M_1, \dots, M_{s-1})$ are fixed. We will reduce the inductive step to two auxiliary claims, which we will then prove separately. To do this, we sample the column $M_s$ by setting all of its entries below the diagonal to be an independent Gaussian $\cN(0, \frac{1}{d})$. This is sufficient to determine whether or not the vertex $s$ is connected to the vertices $s+1, \dots, r$. Let us begin by estimating this probability in our first auxiliary claim.

\begin{claim}\label{claim:auxiliary_1}
Given $M_1, \dots, M_{s-1}$, the probability (over the random choice of $M_s$) that $s$ is connected to all of $s+1, \dots, r$ is at most
{\smaller[0.3]
\begin{equation}\label{eq:auxiliary_1_blue}
    \Pb\bigg[\bigwedge_{s<i\leq r}E_{s, i}\bigg|M[s\!-\!1]\bigg]\leq (1\!-\!p)^{r-s}\exp\bigg(\frac{a\sqrt{d}}{1\!-\!p}\sum_{i=s+1}^r \langle \pi_{s-1}(\by_i), \pi_{s-1}(\by_s)\rangle\!+\!O\big((r\!-\!s)\delta\big)\bigg).
\end{equation}
}
Also, the probability that $s$ is not connected to any of $s+1, \dots, r$ is at most
{\smaller[0.3]
\begin{equation}\label{eq:auxiliary_1_red}
    \Pb\bigg[\bigwedge_{s<i\leq r}\overline{E}_{s, i}\bigg|M[s\!-\!1]\bigg]\leq
    p^{r-s}\exp\bigg(\!-\!\frac{a\sqrt{d}}{p}\sum_{i=s+1}^r \langle \pi_{s-1}(\by_i), \pi_{s-1}(\by_s)\rangle\!+\!O\big((r\!-\!s)\delta\big)\bigg).
\end{equation}
}
\end{claim}

Although the proof of Claim~\ref{claim:auxiliary_1} is quite simple, we postpone it for later. For each outcome of $M_s$, the induction hypothesis provides a bound on the probability of $C_s\wedge B_r$ over the random choice of $M_{s+1}, \dots, M_r$. Thus, in order to perform an induction step, we will integrate this bound over all outcomes of $M_s$ satisfying the event $C_s\wedge B_r$. At this point, it is important to observe that $C_{s-1}=C_{s}\wedge \bigwedge_{s<i\leq r}E_{si}$, i.e. that vertices $s, \dots, r$ form a clique precisely when vertices $s+1, \dots, r$ form a clique and $s$ is connected to all vertices among $s+1, \dots, r$. Hence, by the law of total probability, we can write
{\smaller[0.3]
\begin{align}\label{eq:total_probability_blue}
\Pb\Big[C_{s\!-\!1}\wedge B_r&\Big| M[s\!-\!1]\Big] =\bE_{M_s}\bigg[\Pb\Big[C_{s}\wedge B_r\Big| M_{s}\Big]\bigg| \bigwedge_{s<i\leq r}E_{s, i}, M[s\!-\!1]\bigg]\Pb\bigg[\bigwedge_{s<i\leq r}E_{s, i}\bigg|M[s\!-\!1]\bigg].
\end{align}
}

Since we have already estimated the latter term, we focus on the first one. We use the induction hypothesis to upper bound $\Pb\Big[C_{s}\wedge B_r\Big| M_{s}\Big]$ as follows
{\smaller[0.3]
\begin{equation}\label{eq:induction_hypothesis_blue}
    \Pb\Big[C_{s}\wedge B_r\Big| M_{s}\Big]\leq (1\!-\!p)^{\binom{r-s}{2}}\exp\bigg(\frac{a\sqrt{d}}{1\!-\!p}\!\!\sum_{s<i< j\leq r}\!\!\!\!\langle \pi_s(\by_i), \pi_s(\by_j)\rangle + \frac{a^3}{(1\!-\!p)^3\sqrt{d}}\binom{r\!-\!s}{3}\!+\!O\Big(\frac{(r\!-\!s)^4}{d}\Big)\bigg).
\end{equation}
}

Note that most terms on the right-hand side are independent of the entries of $M_s$. In fact, the only terms which depend on it are $\exp\Big(\frac{a\sqrt{d}}{1-p}\sum_{s<i<j\leq r} \by_i(s)\by_j(s)\Big)$. Thus, we focus on bounding their expectation.

\begin{claim}\label{claim:auxiliary_2}
If the entries of $M_s$ are sampled as Gaussians $\cN(0, \frac{1}{d})$, then we have
{\smaller[0.3]
\begin{align}
    \bE_{M_s}\bigg[\exp\bigg(\frac{a\sqrt{d}}{1\!-\!p}\!\!\sum_{s<i<j\leq r}\!\! \by_i(s)\by_j(s)\bigg)\bigg| \bigwedge_{s<i\leq r}\!\! E_{s, i}, M[s\!-\!1]\bigg]\leq  \exp\bigg(\frac{a^3}{(1\!-\!p)^3\sqrt{d}}\binom{r\!-\!s}{2}\!+\!O\Big(\frac{(r\!-\!s)^3}{d}\!+\!\frac{r-s}{\sqrt{d}}\Big)\bigg)\label{eq:auxiliary_2_blue}    \\
    \bE_{M_s}\bigg[\exp\bigg(\!-\!\frac{a\sqrt{d}}{p}\!\!\sum_{s<i<j\leq r}\!\! \by_i(s)\by_j(s)\bigg)\bigg| \bigwedge_{s<i\leq r}\!\! \overline{E}_{s, i}, M[s\!-\!1]\bigg]\leq \exp\bigg(\!-\!\frac{a^3}{p^3\sqrt{d}}\binom{r\!-\!s}{2}\!+\!O\Big(\frac{(r\!-\!s)^3}{d}\!+\!\frac{r-s}{\sqrt{d}}\Big)\bigg)\label{eq:auxiliary_2_red}
\end{align}
}
\end{claim}

To complete the induction step using the two auxiliary claims, we perform some simple but tedious bookkeeping. By plugging in the bound (\ref{eq:auxiliary_2_blue}) into (\ref{eq:induction_hypothesis_blue}), we obtain
{\smaller[0.3]
\begin{align*}
&\bE_{M_s}\Big[\Pb\big[C_{s}\wedge B_r\big| M_{s}\big]\Big| \bigwedge_{s<i\leq r}E_{s, i}, M[s\!-\!1]\Big]\leq \\
&\leq (1\!-\!p)^{\binom{r-s}{2}}\exp\bigg(\frac{a\sqrt{d}}{1\!-\!p}\!\!\sum_{s<i< j\leq r}\!\!\!\!\langle \pi_{s-1}(\by_i), \pi_{s-1}(\by_j)\rangle\!+\!\frac{a^3}{(1\!-\!p)^3\sqrt{d}}\binom{r\!-\!s}{3}\!+\!O\Big(\frac{(r\!-\!s)^4}{d}\!+\!\frac{(r\!-\!s)^2}{\sqrt{d}}\Big)\bigg) \\
&\hspace{5cm}\cdot \exp\bigg(\frac{a^3}{(1\!-\!p)^3\sqrt{d}}\binom{r\!-\!s}{2}\!+\!O\Big(\frac{ (r\!-\!s)^3}{d}\!+\!\frac{r-s}{\sqrt{d}}\Big)\bigg)\\
&\leq (1\!-\!p)^{\binom{r-s}{2}}\exp\bigg(\frac{a\sqrt{d}}{1\!-\!p}\!\!\sum_{s<i< j\leq r}\!\!\!\!\langle \pi_{s-1}(\by_i), \pi_{s-1}(\by_j)\rangle\!+\!\frac{a^3}{(1\!-\!p)^3\sqrt{d}}\binom{r\!-\!s\!+\!1}{3}\!+\!O\Big(\frac{(r\!-\!s\!+\!1)^4}{d}\!+\!\frac{(r\!-\!s\!+\!1)^2}{\sqrt{d}}\Big)\bigg).
\end{align*}
}

By plugging into the equation (\ref{eq:total_probability_blue}) the bound on $\Pb\Big[\bigwedge_{s<i\leq r}E_{s, i}\Big|M[s-1]\Big]$ from (\ref{eq:auxiliary_1_blue}) and the bound of (\ref{eq:induction_hypothesis_blue}), we obtain
{\smaller[0.3]
\begin{align*}
\Pb\Big[C_{s\!-\!1}\wedge B_r\Big| M[s\!-\!1]\Big]\leq    &(1\!-\!p)^{\binom{r-s}{2}}\exp\bigg(\frac{a\sqrt{d}}{1\!-\!p}\!\!\sum_{s<i< j\leq r}\!\!\!\!\langle \pi_{s-1}(\by_i), \pi_{s-1}(\by_j)\rangle\!+\!\frac{a^3}{(1\!-\!p)^3\sqrt{d}}\binom{r\!-\!s\!+\!1}{3}\\
&\hspace{7cm}\!+\!O\Big(\frac{(r\!-\!s\!+\!1)^4}{d}+\frac{(r\!-\!s\!+\!1)^2}{\sqrt{d}}\Big)\bigg)\\
&\hspace{3cm}\cdot(1\!-\!p)^{r-s}\exp\bigg(\frac{a\sqrt{d}}{1\!-\!p}\sum_{i=s+1}^r \langle \pi_{s-1}(\by_i), \pi_{s-1}(\by_s)\rangle\!+\!O\big((r\!-\!s)\delta\big)\bigg)\\
&\hspace{-4cm}\leq (1\!-\!p)^{\binom{r-s+1}{2}}\exp\bigg(\frac{a\sqrt{d}}{1\!-\!p}\!\!\sum_{s\leq i< j\leq r}\!\!\!\!\langle \pi_{s-1}(\by_i), \pi_{s-1}(\by_j)\rangle\!+\!\frac{a^3}{(1\!-\!p)^3\sqrt{d}}\binom{r\!-\!s\!+\!1}{3}\!+\!O\Big(\frac{(r\!-\!s\!+\!1)^4}{d}+\frac{(r\!-\!s\!+\!1)^2}{\sqrt{d}}\Big)\bigg).
\end{align*}
}
A quick inspection shows that this is precisely the bound we need on $\Pb\big[C_{s-1}\wedge B_r\big|M[s-1]\big]$, and thus our induction step is complete. Let us conclude by noting that nothing in this bookkeeping relied on the fact that we are estimating the probability of $C_s\wedge B_r$. In fact, an analogous proof would apply unchanged to estimate the probability of $I_s\wedge B_r$, and therefore we choose not to spell it out explicitly. Let us now finish the discussion of Proposition~\ref{prop:induction_1} by proving the auxiliary claims.

\begin{proof}[Proof of Claim~\ref{claim:auxiliary_1}.]
Since the proofs of both bounds are analogous, we only show the first one. Recall that $E_{s, i}$ holds whenever $\by_i(s)\by_s(s)\geq -c_p/\sqrt{d}-\langle \pi_{s-1}(\by_i), \pi_{s-1}(\by_s)\rangle$. Note that the only random variable in this equation is $\by_i(s)$, since $\by_s(s)$ was fixed and all of the other ones are determined by $M[s-1]$. Since the variables $\by_{s+1}(s), \dots, \by_{r}(s)$ are independent, so are the events $E_{s, s+1}, \dots, E_{s, r}$.

Thus, it is sufficient to compute the probabilities of the individual events $E_{s, i}$. Thus,
{\smaller[0.3]
\[\Pb\big[E_{s, i}|M[s\!-\!1]\big]=\Pb\bigg[\by_i(s)\geq -\frac{1}{\by_s(s)}\Big(c_p/\sqrt{d}\!+\!\langle \pi_{s-\!1}(\by_i), \pi_{s-\!1}(\by_s)\rangle\Big)\bigg| M[s\!-\!1]\bigg].\]
}
Recall that $\by_s(s)\in (1-2\delta, 1+\delta)$, and note that by Cauchy-Schwarz inequality
$|\langle \pi_{s-1}(\by_i), \pi_{s-1}(\by_s)\rangle|\leq \|\pi_{s-1}(\by_i)\|\cdot \|\pi_{s-1}(\by_s)\|\leq \frac{\alpha^2 \ell}{d}=O(\frac{1}{D\sqrt{d}})$. Therefore the cutoff for $E_{s, i}| M[s\!-\!1]$ to hold equals 
$$b_i=-c_p/\sqrt{d}-\langle \pi_{s-1}(\by_i), \pi_{s-1}(\by_s)\rangle+O(\delta/{\sqrt{d}}).$$ 

Since $\by_i(s)\sim \cN(0, \frac{1}{d})$, log-concavity of $\Phi$ (see (\ref{eqn:log_concavity})) implies that
{\smaller[0.3]
\[\Pb[\by_i(s)\geq b_i]=\Pb[\cN(0, 1)\leq -\sqrt{d}b_i]\leq \Pb[\cN(0, 1)\leq c_p]\exp\Big(\frac{\phi(c_p)}{\Phi(c_p)}(-\sqrt{d}b_i\!-\!c_p)\Big).\]
}
In our setup, we have $\Pb[\cN(0, 1)\leq c_p]=\Phi(c_p)=1-p$ and $\phi(c_p)=a$. Thus, we have
{\smaller[0.3]
\begin{align*}
&\Pb[E_{s, i}]=\Pb[\by_i(s)\geq b_i]\leq (1\!-\!p)\exp\Big(\frac{a\sqrt{d}}{1\!-\!p}\langle \pi_{s-1}(\by_i), \pi_{s-1}(\by_s)\rangle\!+\!O(\delta)\Big),
\end{align*}
}
where we have used that $-\sqrt{d}b_i-c_p=\sqrt{d}\langle \pi_{s-1}(\by_i), \pi_{s-1}(\by_s)\rangle+O(\delta)$. Multiplying over all $s< i\leq r$, we obtain
{\smaller[0.3]
\[\prod_{s<i\leq r}\Pb\Big[E_{s, i}\Big|M[s\!-\!1]\Big]\leq (1\!-\!p)^{r-s}\exp\bigg(\frac{a\sqrt{d}}{1\!-\!p}\sum_{i=s+1}^r \langle \pi_{s-1}(\by_i), \pi_{s-1}(\by_s)\rangle\!+\!O\big((r\!-\!s)\delta\big)\bigg),\]
}
which is sufficient to complete the proof due to the independence of the events $E_{s, i}$.
\end{proof}

\begin{proof}[Proof of Claim~\ref{claim:auxiliary_2}.]
Again, we will focus on proving (\ref{eq:auxiliary_2_blue}), since the proof of (\ref{eq:auxiliary_2_red}) is analogous. Recall that $E_{s, i}$ holds whenever $\by_i(s)\by_s(s)\geq -c_p/\sqrt{d}-\langle \pi_{s-1}(\by_i), \pi_{s-1}(\by_s)\rangle$. Since the diagonal entries $\by_s(s)$ are fixed, the events $E_{s, i}$ are independent and depend only on $\by_i(s)$. So, conditioning on their intersection leaves the variables $\by_{s+1}(s), \dots, \by_{r}(s)$ independent. Moreover, conditioning on the event $E_{s, i}$, each $\by_i(s)$ follows the distribution of a lower truncated Gaussian with the cutoff $b_i=(-c_p/\sqrt{d}-\langle \pi_{s-1}(\by_i), \pi_{s-1}(\by_s)\rangle)/\by_s(s)$.

As we have observed already, truncated Gaussian random variables have variance proxy at most $1/d$, and so we can apply Lemma~\ref{lemma:exponential_moment_special_case} with $\lambda=\frac{a\sqrt{d}}{1-p}$, $k=r-s$, random variables $\by_{s+1}(s), \dots, \by_{r}(s)$ and $S=\sum_{s<i<j\leq r}\by_i(s)\by_j(s)$ to obtain
{\smaller[0.3]
\[\bE[e^{\lambda S}]\leq \exp\bigg( \lambda \bE[S] \!+\! \frac{\lambda^2k^2}{d}\sum_{j=s+1}^r (\bE \by_j(s))^2 \!+\!\frac{4\lambda k}{d}\bigg).\]
}
Here, it is important to verify that $d\geq 4\lambda k$, which follows since $d=D^2\ell^2, \lambda=aD\ell/(1-p)$ and $k=C\ell$, reducing the inequality to $D\ell\geq aC\ell/(1-p)$, which holds since $D$ is sufficiently larger than $C$. Hence, in order to complete the proof, we only need to verify that
{\smaller[0.3]
\[\lambda \bE[S] \!+\! \frac{\lambda^2k^2}{d}\sum_{j=s+1}^r (\bE \by_j(s))^2 \!+\!\frac{4\lambda k}{d}\leq \frac{a^3}{(1\!-\!p)^3\sqrt{d}}\binom{r\!-\!s}{2}\!+\!O\Big(\frac{(r\!-\!s)^2\delta}{\sqrt{d}}\!+\!\frac{(r\!-\!s)^3}{d}\!+\!\frac{r-s}{\sqrt{d}}\Big).\]
}
The last term on the left hand side is bounded by $4\lambda k/d\leq O(\sqrt{d}(r-s)/d)$.

Using the independence of $\by_{s+1}(s), \dots, \by_r(s)$, we can upper bound $\bE[S]$ as follows
{\smaller[0.3]
\[\bE[S]=\sum_{s<i<j\leq r}\bE[\by_i(s)]\bE[\by_j(s)]\leq \frac{1}{2}\Big(\sum_{s<i\leq r}\bE[\by_i(s)]\Big)^2\leq \frac{r\!-\!s}{2}\sum_{s<i\leq r}\bE[\by_i(s)]^2,\]
}
where we have used the inequality between the arithmetic and the quadratic mean in the last step. Thus, we have
{\smaller[0.3]
\[\lambda \bE[S] \!+\! \frac{\lambda^2k^2}{d}\sum_{j=s+1}^r (\bE \by_j(s))^2\leq \lambda \frac{r\!-\!s}{2}\Big(1\!+\!\frac{2\lambda k}{d}\Big)\sum_{i=s+1}^r \bE[\by_i(s)]^2.\]
}

We now need to calculate the expectations of $\by_i(s)$, which we do using Lemma~\ref{lemma:truncated_expectation}. Let us recall that $\by_i(s)$ is a lower truncated Gaussian with the cutoff $b_i=-(c_p/\sqrt{d}+\langle \pi_{s-1}(\by_i), \pi_{s-1}(\by_s)\rangle)/\by_s(s)$. If we write $b_i=-c_p/\sqrt{d}+\eps_i$, then $|\eps_i|=|b_i+c_p/\sqrt{d}|=O(\frac{1}{D\sqrt{d}}+\frac{\delta}{\sqrt{d}})=O(\frac{1}{D\sqrt{d}})$. The last inequality holds since $|\langle \pi_{s-1}(\by_i), \pi_{s-1}(\by_j)\rangle|\leq \frac{\alpha^2\ell}{d}\leq O(\frac{1}{D\sqrt{d}})$ and since $\by_s(s)$ is fixed and in $(1-2\delta, 1+\delta)$. Then, Lemma~\ref{lemma:truncated_expectation} gives
{\smaller[0.3]
\[\bE[\by_i(s)|E_{s, i}]=\frac{\phi(c_p)}{\Phi(c_p)\sqrt{d}}\!+\!O(|\eps_i|)=\frac{a}{(1\!-\!p)\sqrt{d}}\!+\!O\Big(\frac{1}{D\sqrt{d}}\Big).\]
}
Recalling that $\lambda=\frac{a\sqrt{d}}{1-p}$ and $\frac{r-s}{\sqrt{d}}\leq \frac{C\ell}{D\ell}\leq O(\frac{1}{D})$, we get
{\smaller[0.3]
\begin{align*}
\lambda \frac{r\!-\!s}{2} \Big(1\!+\!\frac{2\lambda (r\!-\!s)}{d}\Big)\sum_{i=s+1}^r \bE[\by_i(s)]^2&\leq \frac{a\sqrt{d}(r\!-\!s)}{2(1\!-\!p)}\Big(1\!+\!O\Big(\frac{r\!-\!s}{\sqrt{d}}\Big)\Big)\sum_{i=s+1}^r \left[\frac{a^2}{(1\!-\!p)^2d}\!+\!O\left(\frac{1}{Dd}\right)\right]\\
&\leq \frac{a^3}{(1\!-\!p)^3\sqrt{d}}\binom{r\!-\!s}{2}\Big(1\!+\!O\big(\frac{1}{D}\big)\Big),
\end{align*}
}
which is what we needed to show.
\end{proof}
\end{proof}

\end{document}
